\chapter{Servizi dei sottosistemi}
\label{cap:sdd-servizi}

Questo capitolo definisce l'interfaccia di ciascun sottosistema su due
livelli: le risorse HTTP che espone verso il client e le operazioni che il
suo servizio offre al resto del server. Il primo livello è il contratto
esterno del sistema; il secondo è il contratto fra i livelli architetturali.

\begin{nota}
\textbf{Convenzioni delle risposte.} Le operazioni di lettura rispondono
\id{200}, le creazioni \id{201}, le operazioni che non restituiscono
contenuto \id{204}. Gli errori seguono la gerarchia di
\cref{sec:sdd-errori}: \id{400} dati non validi, \id{401} sessione assente,
\id{403} permessi insufficienti, \id{404} risorsa inesistente,
\id{409} conflitto con lo stato attuale.
\end{nota}

\section{Autenticazione}

\begin{longtable}{@{}L{5.2cm} L{2.6cm} L{6.2cm}@{}}
\toprule
\textbf{Risorsa} & \textbf{Accesso} & \textbf{Descrizione} \\
\midrule
\endfirsthead
\toprule
\textbf{Risorsa} & \textbf{Accesso} & \textbf{Descrizione} \\
\midrule
\endhead
\bottomrule
\endlastfoot

\id{POST /auth/registrazione} & pubblico &
Crea un account con ruolo utente. Non apre la sessione. \\

\id{POST /auth/login} & pubblico &
Verifica le credenziali e deposita il cookie di sessione. \\

\id{POST /auth/logout} & pubblico &
Invalida il cookie di sessione. \\

\id{GET /auth/me} & autenticato &
Profilo dell'utente della sessione corrente. \\

\id{GET /auth/email-disponibile} & pubblico &
Indica se un'email è ancora libera. \\

\end{longtable}

\textbf{Operazioni di \id{AutenticazioneService}}

\begin{longtable}{@{}L{5.8cm} L{8.6cm}@{}}
\toprule
\textbf{Operazione} & \textbf{Contratto} \\
\midrule
\endfirsthead
\toprule
\textbf{Operazione} & \textbf{Contratto} \\
\midrule
\endhead
\bottomrule
\endlastfoot

\id{registra(dati)} &
Crea l'utente con la password sottoposta ad hashing.
\emph{Solleva} \id{ValidationError} se la password non rispetta la politica,
\id{ConflictError} se l'email è già in uso. \\

\id{login(email, password)} &
Restituisce l'utente se le credenziali sono valide.
\emph{Solleva} \id{UnauthorizedError} con lo stesso messaggio sia per email
inesistente sia per password errata. \\

\id{emailEsiste(email)} &
Restituisce un booleano. Non solleva eccezioni. \\

\end{longtable}

\section{Gestione account}

\begin{longtable}{@{}L{5.2cm} L{2.6cm} L{6.2cm}@{}}
\toprule
\textbf{Risorsa} & \textbf{Accesso} & \textbf{Descrizione} \\
\midrule
\endfirsthead
\toprule
\textbf{Risorsa} & \textbf{Accesso} & \textbf{Descrizione} \\
\midrule
\endhead
\bottomrule
\endlastfoot

\id{GET /accounts/me} & autenticato & Profilo dell'utente della sessione. \\
\id{PUT /accounts/me} & autenticato & Aggiorna dati anagrafici e password. \\
\id{DELETE /accounts/me} & autenticato & Elimina il proprio account. \\
\id{GET /accounts/me/badge} & autenticato & Badge conseguiti. \\
\id{GET /accounts} & amministratore & Elenco degli utenti registrati. \\
\id{DELETE /accounts/:id} & amministratore & Elimina l'account indicato. \\

\end{longtable}

\textbf{Operazioni di \id{AccountService}}

\begin{longtable}{@{}L{5.8cm} L{8.6cm}@{}}
\toprule
\textbf{Operazione} & \textbf{Contratto} \\
\midrule
\endfirsthead
\toprule
\textbf{Operazione} & \textbf{Contratto} \\
\midrule
\endhead
\bottomrule
\endlastfoot

\id{getUtente(id)} &
\emph{Solleva} \id{NotFoundError} se l'utente non esiste. \\

\id{getUtenti()} &
Elenco completo, ordinato per cognome e nome. \\

\id{aggiornaProfilo(id, dati)} &
I campi assenti restano invariati. Il cambio password richiede quella
attuale.
\emph{Solleva} \id{NotFoundError}, \id{ConflictError} se la nuova email è già
in uso, \id{UnauthorizedError} se la password attuale non è corretta,
\id{ValidationError} se la nuova password è debole. \\

\id{eliminaUtente(id)} &
Rimuove l'account; feedback, svolgimenti e conseguimenti cadono per vincolo
di integrità. \emph{Solleva} \id{NotFoundError}. \\

\end{longtable}

\section{Gestione tutorial}

\begin{longtable}{@{}L{5.2cm} L{2.6cm} L{6.2cm}@{}}
\toprule
\textbf{Risorsa} & \textbf{Accesso} & \textbf{Descrizione} \\
\midrule
\endfirsthead
\toprule
\textbf{Risorsa} & \textbf{Accesso} & \textbf{Descrizione} \\
\midrule
\endhead
\bottomrule
\endlastfoot

\id{GET /tutorials} & pubblico &
Catalogo, con filtro facoltativo per categoria. \\
\id{GET /tutorials/categorie} & pubblico & Categorie ammesse. \\
\id{GET /tutorials/ricerca} & pubblico & Ricerca per parola chiave. \\
\id{GET /tutorials/:id} & pubblico & Dettaglio di un tutorial. \\
\id{POST /tutorials} & amministratore & Pubblica un tutorial. \\
\id{PUT /tutorials/:id} & amministratore & Aggiorna un tutorial. \\
\id{DELETE /tutorials/:id} & amministratore & Ritira un tutorial. \\
\id{POST /tutorials/immagini} & amministratore & Carica un'immagine di contenuto. \\
\id{DELETE /tutorials/immagini/:nome} & amministratore & Elimina un'immagine di contenuto. \\

\end{longtable}

\textbf{Operazioni di \id{TutorialService}}

\begin{longtable}{@{}L{5.8cm} L{8.6cm}@{}}
\toprule
\textbf{Operazione} & \textbf{Contratto} \\
\midrule
\endfirsthead
\toprule
\textbf{Operazione} & \textbf{Contratto} \\
\midrule
\endhead
\bottomrule
\endlastfoot

\id{getTutorials(categoria?)} &
Senza argomento restituisce l'intero catalogo.
\emph{Solleva} \id{ValidationError} se la categoria non esiste. \\

\id{getTutorial(id)} & \emph{Solleva} \id{NotFoundError}. \\

\id{cercaTutorial(chiave)} &
Chiave vuota restituisce un elenco vuoto senza interrogare il database. \\

\id{creaTutorial(dati)} &
Valida i vincoli di dominio, sanifica il contenuto e persiste.
\emph{Solleva} \id{ValidationError}. \\

\id{aggiornaTutorial(id, dati)} &
Come sopra. \emph{Solleva} anche \id{NotFoundError}. \\

\id{eliminaTutorial(id)} &
Quiz, domande, risposte, svolgimenti e feedback cadono in cascata.
\emph{Solleva} \id{NotFoundError}. \\

\end{longtable}

\section{Gestione quiz}

\begin{longtable}{@{}L{5.2cm} L{2.6cm} L{6.2cm}@{}}
\toprule
\textbf{Risorsa} & \textbf{Accesso} & \textbf{Descrizione} \\
\midrule
\endfirsthead
\toprule
\textbf{Risorsa} & \textbf{Accesso} & \textbf{Descrizione} \\
\midrule
\endhead
\bottomrule
\endlastfoot

\id{GET /quiz/tutorial/:tutorialId} & pubblico &
Quiz del tutorial, privo delle soluzioni. \\
\id{POST /quiz/tutorial/:tutorialId} & amministratore & Definisce il quiz. \\
\id{PUT /quiz/:id} & amministratore & Sostituisce le domande. \\
\id{DELETE /quiz/:id} & amministratore & Rimuove il quiz. \\
\id{POST /quiz/:id/svolgimenti} & autenticato &
Consegna le risposte e ottiene la correzione. \\

\end{longtable}

\textbf{Operazioni di \id{QuizService}}

\begin{longtable}{@{}L{5.8cm} L{8.6cm}@{}}
\toprule
\textbf{Operazione} & \textbf{Contratto} \\
\midrule
\endfirsthead
\toprule
\textbf{Operazione} & \textbf{Contratto} \\
\midrule
\endhead
\bottomrule
\endlastfoot

\id{getQuizPerTutorial(tutorialId)} &
\emph{Solleva} \id{NotFoundError} se il tutorial non ha un quiz. \\

\id{creaQuiz(tutorialId, domande)} &
Inserisce quiz, domande e risposte in un'unica transazione.
\emph{Solleva} \id{NotFoundError} se il tutorial non esiste,
\id{ConflictError} se ha già un quiz, \id{ValidationError} se una domanda non
ha esattamente una risposta corretta o viola i limiti di lunghezza. \\

\id{aggiornaQuiz(quizId, domande)} &
Sostituisce integralmente le domande; gli svolgimenti sono conservati.
\emph{Solleva} \id{NotFoundError}, \id{ValidationError}. \\

\id{eliminaQuiz(quizId)} & \emph{Solleva} \id{NotFoundError}. \\

\id{eseguiQuiz(quizId, utenteId, risposte)} &
Corregge, registra lo svolgimento, ricalcola i quiz superati e delega
l'assegnazione dei badge. Non declassa un quiz già superato.
\emph{Solleva} \id{NotFoundError} se quiz o utente non esistono,
\id{ValidationError} se il quiz non ha domande. \\

\end{longtable}

\section{Gestione feedback}

\begin{longtable}{@{}L{6.0cm} L{2.6cm} L{5.4cm}@{}}
\toprule
\textbf{Risorsa} & \textbf{Accesso} & \textbf{Descrizione} \\
\midrule
\endfirsthead
\toprule
\textbf{Risorsa} & \textbf{Accesso} & \textbf{Descrizione} \\
\midrule
\endhead
\bottomrule
\endlastfoot

\id{GET /feedback/tutorial/:tutorialId} & pubblico & Feedback di un tutorial. \\
\id{GET /feedback/me} & autenticato & Feedback dell'utente. \\
\id{POST /feedback/tutorial/:tutorialId} & autenticato & Inserisce il proprio feedback. \\
\id{PUT /feedback/tutorial/:tutorialId} & autenticato & Modifica il proprio feedback. \\
\id{DELETE /feedback/tutorial/:tutorialId} & autenticato & Elimina il proprio feedback. \\
\id{DELETE /feedback/tutorial/:tutorialId/} \id{utente/:utenteId} & autenticato &
Elimina il feedback indicato; consentito solo se ne è l'autore o è
amministratore. \\

\end{longtable}

\textbf{Operazioni di \id{FeedbackService}}

\begin{longtable}{@{}L{5.8cm} L{8.6cm}@{}}
\toprule
\textbf{Operazione} & \textbf{Contratto} \\
\midrule
\endfirsthead
\toprule
\textbf{Operazione} & \textbf{Contratto} \\
\midrule
\endhead
\bottomrule
\endlastfoot

\id{getFeedbackTutorial(tutorialId)} & Elenco ordinato per data decrescente. \\
\id{getFeedbackUtente(utenteId)} & Elenco dei feedback dell'utente. \\

\id{creaFeedback(utenteId, tutorialId, dati)} &
\emph{Solleva} \id{ValidationError} sui vincoli di valutazione e commento,
\id{NotFoundError} se il tutorial non esiste, \id{ConflictError} se l'utente
ha già valutato quel tutorial. \\

\id{aggiornaFeedback(utenteId, tutorialId, dati)} &
\emph{Solleva} \id{ValidationError}, \id{NotFoundError}. \\

\id{eliminaFeedback(richiedente, autore, tutorialId, isAdmin)} &
Un utente che tenti di eliminare un feedback altrui riceve
\id{NotFoundError}, non \id{ForbiddenError}: negare l'esistenza è preferibile
a confermarla a chi non ha titolo. \\

\end{longtable}

\section{Gestione obiettivi}

\begin{longtable}{@{}L{5.2cm} L{2.6cm} L{6.2cm}@{}}
\toprule
\textbf{Risorsa} & \textbf{Accesso} & \textbf{Descrizione} \\
\midrule
\endfirsthead
\toprule
\textbf{Risorsa} & \textbf{Accesso} & \textbf{Descrizione} \\
\midrule
\endhead
\bottomrule
\endlastfoot

\id{GET /obiettivi} & pubblico & Obiettivi raggiungibili e relative soglie. \\
\id{POST /obiettivi} & amministratore & Definisce un obiettivo. \\
\id{PUT /obiettivi/:nome} & amministratore & Aggiorna un obiettivo. \\
\id{DELETE /obiettivi/:nome} & amministratore & Elimina un obiettivo. \\

\end{longtable}

\textbf{Operazioni di \id{ObiettivoService}}

\begin{longtable}{@{}L{5.4cm} L{9.0cm}@{}}
\toprule
\textbf{Operazione} & \textbf{Contratto} \\
\midrule
\endfirsthead
\toprule
\textbf{Operazione} & \textbf{Contratto} \\
\midrule
\endhead
\bottomrule
\endlastfoot

\id{getObiettivi()} & Catalogo ordinato per soglia crescente. \\

\id{getBadgeUtente(utenteId)} &
Conseguimenti abbinati al rispettivo obiettivo; i conseguimenti il cui
obiettivo è stato eliminato sono omessi. \\

\id{valutaConseguimenti(utenteId, quizSuperati, conn?)} &
Assegna gli obiettivi la cui soglia è raggiunta e non ancora conseguiti;
restituisce quelli sbloccati dalla chiamata. È idempotente: invocata due
volte con lo stesso stato non assegna nulla la seconda volta. Accetta una
connessione transazionale. \\

\id{creaObiettivo(dati)} &
\emph{Solleva} \id{ValidationError}, \id{ConflictError} se il nome è già usato. \\

\id{aggiornaObiettivo(nome, dati)} &
\emph{Solleva} \id{NotFoundError}, \id{ValidationError}. \\

\id{eliminaObiettivo(nome)} &
I conseguimenti associati cadono in cascata.
\emph{Solleva} \id{NotFoundError}. \\

\end{longtable}
